% Pursuing Stacks (cotes 134-2 to 134-8): archive cote 134-7, site page 57 (PDF page 58), typescript page 499.
% typescript. Draft compared: ps6.tex:2817-2869.
% AI-transcribed by Claude Opus 5.5 (claude-opus-5-5, adaptive thinking, effort high), Message Batches API, batch ps4 (msgbatch_012Aoc1NuskATPsyBeKSYK7J), 2026-11-12; not yet checked by eye.
\scanpage{134-7/57}
\mspage{499}

\uline{the} basic building block for building up the most general semisimple
groups,  whereas $P^{1}_{\hand{\underline{C}}}$ may be viewed as being the most significant homoge-%
neous space under this group,  namely the first and most elementary case
of flag manifolds.  In view of this significance of $S^2$ and $S^3$,  it is all
the more remarkable that no simple, non plethoric semisimplicial model
say is available,  in terms say of a semisimplicial group having the
homotopy type of \xout{\unsure{z}}
\[ S^3 \hand{\sim} S\unsure{l}(2,\underline{C} )= G \hand{,} \]
and a subgroup having the homotopy type of
\[ S^1 \hand{\sim} \underline{C}^{\hand{*}} = G\unsure{l}(\unsure{1},\underline{C}) \hand{\sim} K(\underline{Z},1) \]
and playing the part of a Borel subgroup or a maximal torus,  in such a
way that the quotient will have the homotopy type of
\[ S^2 = P^1_{\underline{C}} \hand{\simeq} S^3/S^1 \hand{\simeq} G/B \quad . \]
It would be tempting now to try and construct such a model of the situa-%
tion in terms of semisimplicial unipotent bund\over{\illegible{1}}{les} over $\underline{Z}$  - which would
at the same time display the homotopy groups of $S^2$ and $S^3$ (not much of a
difference\hand{!}) via formula (9).  All the more tempting of course,  as it is
felt that the geometric objects and their relationship, \xout{which} the homotopy
shadow of which we want to modelize \deleted{\illegible{1}} schematically,  are themselves\inserted{already,}
basically, \inserted{most beautiful} schemes over Spec($\underline{Z}$)\hand{!}

\margin{! (with a vertical line along the next three lines)}Another way of getting a display of the homotopy groups of $S^2$ and
$S^3$ would be in terms of a (discrete) model of the situation above,  in
terms of ``hemispherical complexes'' rather than semisimplicial ones.  On
the other hand,  there is no reason why a theory of schematic homotopy
types could not be carried through as well,  using hemispherical comple-%
xes rather than semisimplicial ones.  The latter \inserted{kind of complexes} have the advantage \deleted{\unsure{of}}
that they have become \xout{familiar}  thoroughly familiar \xout{\unsure{by}} \inserted[typed]{through} constant use
by topologists and homotopy people for thirty years or so  -  the former
however and newcomers,  have the advantage of still greater formal simpli-%
city (just two boundary operations,  and just one degeneracy),  and more
important still,  of allowing for a direct computational description of
the homotopy invariants,  in the discrete set-up.  When working with
hemispherical complexes of unipotent bundles as models for schematic
homotopy types,  we'll get then \uline{two} highly differen\hand{t} descriptions of the
homotopy invariants $\hand{\pi}_i$ ,  one by the ``infinitesimal'' formula(9)  interpre-%
ting them as derived functors of the Lie functor,  the other one via the

\margin{57}
